\chapter{Matrix Factorisations: SVD, PCA, NMF, Recommenders and t-SNE}
\companion{StatQuest: Principal Component Analysis (PCA), Step by Step}{\ytid{FgakZw6K1QQ}}

The InfoEdge hire was asked how PCA uses SVD, what a scree plot is, and how t-SNE and its perplexity work. Matrix
factorisation is also how classic recommenders (``jobs you may like'', Jeevansathi's collaborative models) work. This chapter
starts from ``what is a matrix'' and ends with recommender maths you can do on paper.

\section{What ``factorising a matrix'' means}
A matrix is a table of numbers. Factorising it means writing it as a product of simpler matrices, just as $12 = 3\times4$.
Why? Because the factors reveal structure and compress data.

\begin{example}[: a rank-1 table]
Ratings of 3 users for 3 jobs:
$A = \begin{pmatrix}2 & 4 & 6\\1 & 2 & 3\\3 & 6 & 9\end{pmatrix} = \begin{pmatrix}2\\1\\3\end{pmatrix}\begin{pmatrix}1 & 2 & 3\end{pmatrix}$.
Nine numbers are explained by six: a ``user enthusiasm'' and a ``job attractiveness''. Every row is a multiple of every other: the
\textbf{rank} is 1 (only one independent pattern).
\end{example}
Real data are never exactly low-rank, but they are often \emph{approximately} low-rank: a few patterns explain most of the variation.
Factorisation finds those patterns.

\section{A quick refresher: vectors, matrices, eigenvectors}
\begin{itemize}
\item Multiplying a matrix by a vector transforms the vector (stretches, rotates, projects).
\item An \textbf{eigenvector} $\mathbf v$ of a square matrix $A$ is a direction that $A$ only stretches: $A\mathbf v = \lambda\mathbf v$; the
  stretch factor $\lambda$ is the \textbf{eigenvalue}.
\item A symmetric matrix (like a covariance matrix) has real eigenvalues and perpendicular eigenvectors.
\item The \textbf{covariance matrix} $\Sigma$ of centred data $X$ ($n\times d$) is $\frac{1}{n-1}X^\top X$: entry $(j,k)$ is the covariance of
  features $j$ and $k$; the diagonal holds variances.
\end{itemize}

\section{SVD: the master factorisation}
\textbf{Every} real $m\times n$ matrix can be written
\[ A = U\Sigma V^\top, \]
\begin{itemize}
\item $V$ ($n\times n$): orthonormal columns $\mathbf v_i$, the \textbf{right singular vectors}: directions in the input (column) space.
\item $\Sigma$: diagonal with \textbf{singular values} $\sigma_1\ge\sigma_2\ge\dots\ge0$: how much $A$ stretches each direction.
\item $U$ ($m\times m$): orthonormal columns $\mathbf u_i$, the \textbf{left singular vectors}: where those directions land.
\end{itemize}
Geometric meaning: any linear map is a rotation ($V^\top$), then a stretch along the axes ($\Sigma$), then another rotation ($U$).
Equivalently $A = \sum_i\sigma_i\mathbf u_i\mathbf v_i^\top$: a sum of rank-1 layers, most important first.

Connections: $A^\top A = V\Sigma^2V^\top$, so the $\mathbf v_i$ are eigenvectors of $A^\top A$ and $\sigma_i^2$ its eigenvalues. Rank $=$
number of non-zero singular values. For a symmetric positive semi-definite matrix, SVD and eigendecomposition coincide. Thin SVD of a
$1000\times50$ full-rank matrix: $U$ is $1000\times50$, $\Sigma$ $50\times50$, $V$ $50\times50$.

\subsection{Low-rank approximation (Eckart--Young)}
Keep only the top $k$ terms: $A_k = \sum_{i=1}^k\sigma_i\mathbf u_i\mathbf v_i^\top$. This is the \emph{best} rank-$k$ approximation of $A$ in
both Frobenius and spectral norms. The squared Frobenius error equals the sum of the discarded $\sigma_i^2$.
\begin{example}[: compression]
Singular values 5, 3, 1. Rank-1 error $=\sqrt{9 + 1} = 3.16$. Rank 2 keeps $(25 + 9)/35 = 97\%$ of the squared norm. A $100\times200$ image
stored with rank 10 needs $10(100 + 200 + 1) = 3{,}010$ numbers instead of 20{,}000.
\end{example}

\section{PCA is SVD of the centred data}
\subsection{The goal}
Find a few new axes (\textbf{principal components}) that capture as much of the data's variation as possible, so we can describe each
point with fewer numbers while losing little information.

\subsection{Step by step}
\begin{enumerate}
\item \textbf{Centre} each feature (subtract its mean). Usually also \textbf{standardise} (divide by its standard deviation), otherwise
  features with large units dominate.
\item Compute the covariance matrix $\Sigma = \frac{1}{n-1}X^\top X$.
\item Find its eigenvectors (the principal directions) and eigenvalues (the variance along each).
\item Sort by eigenvalue; keep the top $k$. Project: $Z = XV_k$ (each point's coordinates on the new axes).
\end{enumerate}
\textbf{Via SVD}: $X = U\Sigma V^\top$ directly gives the directions $V$; the variances are $\sigma_i^2/(n-1)$; the projected coordinates
are $U\Sigma$. Libraries use SVD because it avoids forming $X^\top X$ (which squares the condition number and loses precision).

\begin{example}[: PCA on two features by hand]
Covariance matrix $\begin{pmatrix}2 & 1\\1 & 2\end{pmatrix}$ (two standardised-ish features with covariance 1).
Eigenvalues: $\det\begin{pmatrix}2-\lambda & 1\\1 & 2-\lambda\end{pmatrix} = (2-\lambda)^2 - 1 = 0\Rightarrow\lambda = 3, 1$.
Eigenvectors: $(1,1)/\sqrt2$ for 3 and $(1,-1)/\sqrt2$ for 1. PC1 is ``the sum of the two features'' and carries $3/4 = 75\%$ of the variance.
\end{example}

\subsection{How many components? The scree plot}
A \textbf{scree plot} shows eigenvalues (or explained-variance ratio $\lambda_i/\sum\lambda$) against component number. Look for the
``elbow'' after which the curve flattens (``scree'' is the rubble at the foot of a cliff). Alternatives: keep enough components for
90--95\% cumulative variance; Kaiser rule (eigenvalue $>1$ on standardised data); or choose $k$ by cross-validated downstream performance.
\begin{example}[: explained variance]
Eigenvalues 6, 3, 1: PC1 explains 60\%, PC1+PC2 90\%.
\end{example}

\subsection{Limitations of PCA}
Linear only; maximises variance, which need not be what matters for a target (Chapter 10); sensitive to scaling and outliers; components
are mixtures of all features and hard to interpret; assumes the mean and covariance summarise the data.

\section{NMF: non-negative parts}
\textbf{Non-negative matrix factorisation}: $A\approx WH$ with all entries of $W$ and $H$ non-negative. Because nothing can cancel,
factors are additive ``parts''. On a TF-IDF document--term matrix, each row of $H$ is a \emph{topic} (a non-negative combination of words) and
each row of $W$ says how much of each topic a document contains. Fitted by multiplicative updates or alternating least squares; not unique;
needs non-negative input.

\section{Matrix factorisation for recommendations}
\subsection{The setting}
A huge user $\times$ item matrix $R$ (seekers $\times$ jobs) with a few observed entries (ratings, applies, clicks); most are unknown, not zero.
Goal: predict the unknown entries.

\subsection{The model}
Give each user $u$ a vector $\mathbf p_u\in\R^k$ and each item $i$ a vector $\mathbf q_i\in\R^k$ (latent factors, $k\approx 32$--$200$). Predict
\[ \hat r_{ui} = \mu + b_u + b_i + \mathbf p_u^\top\mathbf q_i, \]
with global mean $\mu$, user bias $b_u$ (this user rates/applies a lot) and item bias $b_i$ (popular job). Factors are learned so that
dot products reproduce the observed entries:
\[ \min\sum_{(u,i)\ \text{observed}}(r_{ui} - \hat r_{ui})^2 + \lambda(\norm{\mathbf p_u}^2 + \norm{\mathbf q_i}^2 + b_u^2 + b_i^2). \]
Latent dimensions often become interpretable (``remote-friendly'', ``startup vs MNC'', ``data roles'').

\begin{example}[: predicting a rating]
$\mu = 3.5$, $b_u = 0.2$, $b_i = -0.3$, $\mathbf p_u = (1, 2)$, $\mathbf q_i = (3, 0.5)$: $\mathbf p_u^\top\mathbf q_i = 3 + 1 = 4$; $\hat r = 3.5 + 0.2 - 0.3 + 4 = 7.4$
(on an unbounded scale; clip if needed).
\end{example}

\subsection{Why not plain SVD?}
SVD needs a complete matrix; filling unknowns with 0 says ``the user dislikes every unseen job'', which is false. Matrix factorisation only fits
the observed entries.

\subsection{Training: SGD and ALS}
\textbf{SGD}: for each observed $(u,i)$, error $e = r - \hat r$; update $\mathbf p_u\leftarrow\mathbf p_u + \eta(e\,\mathbf q_i - \lambda\mathbf p_u)$ and
$\mathbf q_i\leftarrow\mathbf q_i + \eta(e\,\mathbf p_u - \lambda\mathbf q_i)$.
\begin{example}[: one SGD step]
$r = 4$, $\mathbf p = (1,0)$, $\mathbf q = (2,1)$, $\eta = 0.1$, $\lambda = 0$: $\hat r = 2$, $e = 2$. $\mathbf p\leftarrow(1,0) + 0.2(2,1) = (1.4, 0.2)$;
$\mathbf q\leftarrow(2,1) + 0.2(1,0) = (2.2, 1)$ (using the old $\mathbf p$).
\end{example}
\textbf{ALS (alternating least squares)}: fix all item vectors and solve each user vector exactly (a small ridge regression), then fix users and
solve items; repeat. Each subproblem is convex and independent per user, so ALS parallelises well (Spark MLlib) and handles implicit feedback.

\subsection{Implicit feedback and cold start}
Clicks and applies are \emph{implicit}: an absent click is not a dislike. Implicit ALS treats every entry with a confidence weight
($1 + \alpha\cdot\text{count}$) and a binary preference. \textbf{Cold start} (new job, new user): no interactions, no factors. Use content
features (title, skills, location embeddings), hybrid models (factorisation machines, two-tower networks), popularity priors, and exploration.

\section{t-SNE and UMAP: seeing clusters}
PCA is linear and preserves large-scale variance. For \emph{visualising} clusters in high-dimensional data (resume embeddings), use
non-linear methods that preserve \emph{neighbourhoods}.

\subsection{t-SNE in steps}
\begin{enumerate}
\item In the original space, turn distances into probabilities that $j$ is $i$'s neighbour: a Gaussian around each point, $p_{j|i}\propto
  \exp(-\norm{\mathbf x_i - \mathbf x_j}^2/2\sigma_i^2)$; symmetrise.
\item Each point's bandwidth $\sigma_i$ is chosen so that its neighbourhood has a set \textbf{perplexity} (roughly the effective number of
  neighbours, typically 5--50): dense regions get small $\sigma_i$, sparse regions large.
\item In 2-D, define similarities with a heavy-tailed \textbf{Student-t} distribution: $q_{ij}\propto(1 + \norm{\mathbf y_i - \mathbf y_j}^2)^{-1}$.
\item Move the 2-D points by gradient descent to minimise $KL(P\,\|\,Q)$.
\end{enumerate}
Why the t-distribution: its heavy tails let moderately distant points sit far apart in 2-D, relieving the ``crowding problem'' and producing
clear, separated clusters.

\subsection{How to read (and not read) a t-SNE plot}
\begin{itemize}
\item Clusters and local neighbourhoods are meaningful.
\item Cluster \emph{sizes} and the \emph{distances between clusters} are not reliably meaningful.
\item Results depend on perplexity and random seed; run several.
\item It is for visualisation: no \code{transform} for new points; do not feed t-SNE coordinates to a classifier.
\item \textbf{Perplexity}: low values focus on very local structure (many small clumps); high values capture more global structure but can
  merge clusters.
\end{itemize}
\textbf{UMAP} is a faster alternative that preserves more global structure and can transform new points.

\begin{practical}[: factorisations at InfoEdge]
\begin{itemize}
\item \textbf{Collaborative models} (Jeevansathi slide: ``content models, activity-based models, collaborative models''): MF/ALS on
  profile-view and interest interactions; ``users who showed interest in similar profiles''.
\item \textbf{Job recommendations} (``Reco'' in the Seeker AI engines): MF or two-tower models on apply data, with content features for cold
  start (most jobs live only weeks, so cold start is the norm).
\item \textbf{LSA} (truncated SVD of TF-IDF) as a quick semantic search baseline; PCA to compress embeddings for storage in a vector index.
\item t-SNE/UMAP of resume embeddings to sanity-check that the taxonomy's clusters (Data Science, DevOps, Sales) separate.
\end{itemize}
\end{practical}

\stuck{StatQuest: t-SNE, Clearly Explained}{\ytid{NEaUSP4YerM}}
\section{Interview questions}

\begin{qa}{\pyq{INT: PCA via SVD and the scree plot} Explain PCA and how it uses SVD. What is a scree plot?}
\oneline{PCA finds orthogonal directions of maximum variance in centred data; the right singular vectors of the centred data matrix are those
directions, and the squared singular values divided by $n-1$ are the variances; a scree plot shows those variances in decreasing order to choose
how many components to keep.}
\why
\textbf{Objective.} Find unit vector $\mathbf v$ maximising the variance of projections $\Var(X\mathbf v) = \mathbf v^\top\Sigma\mathbf v$. With a Lagrange
multiplier for $\norm{\mathbf v} = 1$: $\Sigma\mathbf v = \lambda\mathbf v$, so the best $\mathbf v$ is the top eigenvector and the variance captured is
$\lambda$. The next component maximises variance subject to being orthogonal to the first, and so on.
\textbf{SVD link.} $X = U\Sigma_sV^\top$ gives $X^\top X = V\Sigma_s^2V^\top$: eigenvectors $V$, eigenvalues $\sigma_i^2$. So compute SVD of the
centred $X$ and read off $V$ (loadings), $\sigma_i^2/(n-1)$ (variances) and $U\Sigma_s$ (scores). SVD is numerically better than forming $X^\top X$.
\textbf{Scree plot.} Plot $\lambda_i$ or explained-variance ratio; choose $k$ at the elbow or by a cumulative threshold.
\worked Eigenvalues 6, 3, 1: 60\%, 30\%, 10\%; two components keep 90\%. Centred data with $n = 101$ and $\sigma_1 = 20$: variance along PC1 $= 400/100 = 4$.
\pushes \emph{``Why centre?''} Without centring, the first component points toward the mean rather than along the spread. \emph{``Why standardise?''}
So a feature in rupees does not dominate one in years. \emph{``Is PCA feature selection?''} No, each component mixes all features.
\end{qa}

\begin{qa}{\pyq{INT: t-SNE and perplexity} Why is t-SNE useful for seeing clusters, and what does perplexity control?}
\oneline{t-SNE places points in 2-D so that points that are neighbours in high dimensions remain neighbours, using a heavy-tailed similarity in 2-D
that pulls clusters apart; perplexity sets the effective number of neighbours each point considers, trading local against more global structure.}
\why Perplexity $= 2^{H(P_i)}$, where $H$ is the entropy of point $i$'s neighbour distribution; each point's Gaussian width is tuned so the perplexity
matches the chosen value, so dense and sparse regions are treated fairly. Low perplexity (5): fragmented, very local; high (50--100): smoother,
more global, small clusters may merge.
\pushes \emph{``Can I trust distances between clusters?''} No. \emph{``Why not use it as features?''} It is non-parametric, stochastic, and has no
out-of-sample mapping; use UMAP or PCA for features.
\end{qa}

\begin{qa}{\pyq{INT: PCA vs t-SNE vs LDA} Compare them.}
\oneline{PCA is linear, unsupervised and preserves global variance; LDA is linear, supervised and maximises class separation (at most $K-1$
dimensions); t-SNE is non-linear, unsupervised and preserves local neighbourhoods for visualisation.}
\why Use PCA for compression, denoising and decorrelation before modelling; LDA for supervised reduction or a linear classifier; t-SNE/UMAP for
looking at cluster structure.
\end{qa}

\begin{qa}{What is the SVD and what do $U$, $\Sigma$ and $V$ mean geometrically?}
\oneline{$A = U\Sigma V^\top$: $V^\top$ rotates the input onto special orthogonal directions, $\Sigma$ stretches each by a singular value, and $U$
rotates the result into the output space.}
\why The unit circle maps to an ellipse whose semi-axes have lengths $\sigma_i$ along $\mathbf u_i$. Singular values are always non-negative; the
number of non-zero ones is the rank.
\end{qa}

\begin{qa}{\pyq{reported linear algebra: rank, null space} What is the rank of a matrix and its null space, and how does SVD reveal both?}
\oneline{Rank is the number of linearly independent columns (the dimension of the column space); the null space is the set of vectors the matrix sends
to zero; in the SVD, rank $=$ number of non-zero singular values, and the right singular vectors with zero singular value span the null space.}
\why Rank--nullity: rank $+$ dim(null space) $=$ number of columns. In regression, a non-trivial null space of $X$ means some combination of features
is identically zero (perfect multicollinearity), so coefficients are not unique: adding any null-space vector to $\mathbf w$ leaves predictions unchanged.
\worked $\begin{pmatrix}1 & 2\\2 & 4\end{pmatrix}$: rank 1; null space spanned by $(2, -1)$ because $1\cdot2 + 2\cdot(-1) = 0$.
\end{qa}

\begin{qa}{Explain the Eckart--Young theorem and compute the rank-1 approximation error for singular values 5, 3, 1.}
\oneline{Truncating the SVD to the top $k$ singular values gives the best rank-$k$ approximation; the Frobenius error is $\sqrt{3^2 + 1^2} = 3.16$.}
\end{qa}

\begin{qa}{Find the principal components of the covariance matrix $\begin{pmatrix}2 & 1\\1 & 2\end{pmatrix}$.}
\oneline{Eigenvalues 3 and 1 with eigenvectors $(1,1)/\sqrt2$ and $(1,-1)/\sqrt2$; PC1 explains 75\% of the variance.}
\end{qa}

\begin{qa}{\deep Can PCA hurt a classifier?}
\oneline{Yes: PCA keeps high-variance directions, and the direction that separates the classes may have low variance and be discarded; also it can mix
informative and noisy features.}
\why The parallel-cigars example: the class difference lies along the narrow axis, PC2, which a ``keep 95\% variance'' rule might drop. Choose $k$ by
validation performance, or use supervised reduction (LDA, PLS).
\end{qa}

\begin{qa}{\deep Two features are perfectly correlated ($x_2 = 2x_1$). What does PCA give?}
\oneline{One component with all the variance and a second with eigenvalue exactly zero: the data lie on a line, so the covariance matrix has rank 1.}
\why The zero eigenvector $(2, -1)/\sqrt5$ spans the null space of the centred data: a combination that is always 0.
\end{qa}

\begin{qa}{How does matrix factorisation predict a seeker's interest in a job they have never seen?}
\oneline{It represents each seeker and job as vectors learned from observed interactions so that dot products (plus biases) reproduce known interests;
the dot product of an unseen pair is the prediction.}
\why Similar seekers get similar vectors because they interacted with similar jobs; so a seeker inherits the tastes of their ``neighbours'' in latent
space. Worked: $3.5 + 0.2 - 0.3 + 4 = 7.4$.
\end{qa}

\begin{qa}{Why not run ordinary SVD on a sparse ratings matrix?}
\oneline{SVD requires every entry, and filling unknown ratings with zeros treats ``not seen'' as ``hated'', biasing all factors toward predicting zero;
matrix factorisation fits only observed entries with regularisation.}
\end{qa}

\begin{qa}{ALS vs SGD for matrix factorisation: which and why?}
\oneline{ALS alternates exact least-squares solves for users and items, which is easily parallelised and handles implicit-feedback weighting efficiently;
SGD is simpler, works well for explicit ratings, and handles extra terms easily but needs learning-rate tuning.}
\worked ALS with $k = 1$: a user rated items with $q = 1$ and $q = 2$ as 3 and 4; with $\lambda = 1$: $p = \frac{1\cdot3 + 2\cdot4}{1 + 4 + 1} = 11/6 = 1.83$.
\end{qa}

\begin{qa}{\deep How do you recommend a job posted five minutes ago, with zero interactions?}
\oneline{With content-based signals: embed the job from its title, skills, location and salary (text encoders, taxonomy features), match it to seekers
whose profiles or past applies are similar, and use exploration to collect early feedback; hybrid models map content features into the same latent
space as collaborative factors.}
\why Two-tower models learn a job tower from content features, so new jobs get vectors immediately. Also use popularity of similar jobs as a prior and
bandit-style exploration so new jobs get some impressions.
\end{qa}

\begin{qa}{What is NMF and how does it differ from SVD?}
\oneline{NMF factorises a non-negative matrix into non-negative factors, giving additive, parts-based, interpretable components (like topics); SVD factors
can be negative, are orthogonal and optimal in squared error, but are harder to interpret.}
\end{qa}

\begin{qa}{How many parameters does a factorisation with 10{,}000 users, 5{,}000 jobs and $k = 32$ have (no biases)?}
\oneline{$(10{,}000 + 5{,}000)\times32 = 480{,}000$.}
\end{qa}

\begin{qa}{\deep LoRA fine-tunes a large model by learning $\Delta W = BA$ with rank $r$. How is this a matrix factorisation, and how many parameters does
it save for a $4096\times4096$ weight with $r = 8$?}
\oneline{It constrains the weight update to be low-rank, the product of a $4096\times8$ and an $8\times4096$ matrix: $65{,}536$ parameters instead of
$16{,}777{,}216$, a 256-fold saving.}
\why The hypothesis is that the change needed to adapt a pretrained model has low intrinsic rank. This is exactly a truncated factorisation of the update.
\end{qa}

\begin{qa}{What is the condition number of a matrix and why does it matter in ML?}
\oneline{The ratio of largest to smallest singular value; a large condition number means the problem is ill-conditioned: small changes in data cause
large changes in the solution, and gradient descent converges slowly.}
\worked Singular values 100, 10, 0.1: condition number 1{,}000. Ridge improves it to $(\sigma_1^2 + \lambda)/(\sigma_{min}^2 + \lambda)$ for the normal equations.
\end{qa}

\begin{qa}{\deep Is PCA the same as a linear autoencoder?}
\oneline{A linear autoencoder with a $k$-unit bottleneck trained with squared error learns the same subspace as the top $k$ principal components, though
its weights need not be orthogonal or ordered.}
\why Non-linear autoencoders generalise PCA to curved manifolds (Chapter 22).
\end{qa}

\begin{qa}{What is Latent Semantic Analysis?}
\oneline{Truncated SVD of a document--term (TF-IDF) matrix; documents and words are represented in a low-dimensional ``concept'' space where synonyms that
co-occur with similar words end up close.}
\why Use \code{TruncatedSVD} rather than \code{PCA} because it works on sparse matrices without centring (centring would make them dense).
\end{qa}
